% Generated draft: requires mathematical and source review.
\documentclass[11pt]{article}
\usepackage[margin=25mm]{geometry}
\usepackage{fontspec}
\setmainfont{Latin Modern Roman}
\usepackage{amsmath,amssymb,mathtools,unicode-math}
\setmathfont{Latin Modern Math}
\usepackage{enumitem,needspace,xurl,hyperref}
\setlength{\parindent}{0pt}
\setlength{\parskip}{5pt}
\setlength{\emergencystretch}{3em}
\newcommand{\N}{\mathbb{N}}
\newcommand{\Z}{\mathbb{Z}}
\newcommand{\Q}{\mathbb{Q}}
\newcommand{\R}{\mathbb{R}}
\newcommand{\C}{\mathbb{C}}
\newcommand{\F}{\mathbb{F}}
\newcommand{\A}{\mathbb{A}}
\newcommand{\PP}{\mathbb{P}}
\newcommand{\E}{\mathbb{E}}
\newcommand{\eps}{\varepsilon}
\newcommand{\dd}{\,\mathrm{d}}
\begin{document}
\section*{UnsolvedMath ID: 30000094}
\textbf{Source problem code:} catalogue entry
\subsection{Definitions and mathematical statement}
The mathematical question concerns warped compactly supported wavelet bases and asks which hypotheses preserve shrinkage, unconditionality, and \(p\)-democracy. To express the desired properties, let \((\psi_\lambda)_{\lambda\in\Lambda}\) denote the starting wavelet basis, and write \((w_\lambda)\) for the basis obtained by warping with \(G\). A simple composition convention would be \(w_\lambda(x)=\psi_\lambda(G(x))\); a Jacobian-weighted transformation is a different convention. The source does not give the actual formula, the underlying domain and measure, or the assumptions on \(G\).

For a specified \(L^p\) norm, unconditional shrinkage means that there is \(C<\infty\) such that, for every finitely supported coefficient family \((a_\lambda)\) and scalars \(|\varepsilon_\lambda|\leq1\),
\[
\left\|\sum_\lambda\varepsilon_\lambda a_\lambda w_\lambda\right\|_p
\leq C\left\|\sum_\lambda a_\lambda w_\lambda\right\|_p.
\]
This finite-sum inequality accompanies the requirement that the warped family is a basis of the intended function space. For \(1\leq p<\infty\), a conventional normalized formulation of \(p\)-democracy uses \(e_\lambda=w_\lambda/\|w_\lambda\|_p\) and asks for constants \(0<c\leq C'<\infty\) such that
\[
c|F|^{1/p}\leq\left\|\sum_{\lambda\in F}e_\lambda\right\|_p
\leq C'|F|^{1/p}
\quad\text{for every finite nonempty }F\subset\Lambda.
\]
The constants are independent of coefficients, shrinkage multipliers, and finite sets, as applicable; their allowed dependence on \(G\), \(p\), the measure, and the original wavelets needs specification. The source's precise democracy normalization and admissible \(p\) range must be checked before this conventional display is adopted as its exact definition.

The research task is to identify sufficient or characterizing hypotheses on the warp, original wavelets, exponent, and function space under which these properties survive, together with the approximation properties named in the original wording. ``Sufficiently regular'' does not specify differentiability, invertibility, monotonicity, derivative bounds, or weight conditions, so it is not a complete quantified assumption.

There is a material mismatch in the supplied record: its literature assessment checked 2026-08-21 discusses sparse recovery from random Fourier samples and cites compressed-sensing papers, whereas its statement concerns warped wavelet bases. The assessment says exact recovery is possible under sparsity and structural assumptions and that an unrestricted sampling question has no affirmative answer; those claims do not resolve the warped-wavelet question. The partially solved label is retained as metadata, with this mismatch explicit.
\subsection{Short English statement}
Determine assumptions under which warping a compactly supported wavelet basis preserves unconditional coefficient shrinkage and \(p\)-democracy. The warp and norm conventions are missing, and the supplied compressed-sensing literature assessment addresses a different question.
\subsection{Sources}
\begin{itemize}
\item Local UnsolvedMath registry, permanent ID 30000094, catalogue entry, \emph{Shrinkage and Democracy of Warped Wavelet Bases}.
\item \emph{Wavelet and Multiscale Methods} (2005), Oberwolfach Reports, pp. 1789--1844, DOI \url{https://doi.org/10.4171/owr/2004/34}. The supplied statement-verification field says that the two properties preceding the report's question were restored.
\item Supplied but apparently mismatched literature reference: Cand\`es--Romberg--Tao, \emph{Robust uncertainty principles}, \url{https://arxiv.org/abs/math/0409186}; supplied IEEE published version \url{https://doi.org/10.1109/TIT.2005.862083}.
\item Supplied literature triage dated 2026-08-21 discusses random Fourier reconstruction rather than the warped-basis problem; its relevance requires source reconciliation.
\end{itemize}
\end{document}
